\section{Introduction}
\label{sec:intro}

The public ALMA archive holds thousands of hours of millimetre and
submillimetre observations taken for other reasons, mostly studies of dust
discs and planet formation. Many of those fields contain a nearby star.
This
paper searches them for a narrowband transmitter at the position of such a
star. No such signal was found.

Searches for \emph{technosignatures} -- detectable imprints of technology,
intended or not \citep{Sheikh2020,SETIreview2025} -- have examined many
thousands of stars since the first proposals
\citep{CocconiMorrison1959,Drake1960}, from Project Phoenix
\citep{Tarter2001,Backus2002} to the targeted, commensal and
machine-learning-assisted surveys of today
\citep{Zhang2020,Czech2021,FAST3IATLAS2025,LOFARdualsite2023,Choza2024,Morrison2023,Tremblay2024,SRT2025,Ma2023,GLOBULAR2025,AnomalyParkesGBT2025}.
Even so, the last published accounting of that effort covers only
\HaystackFrac{} of the space of transmitter properties \citep{Wright2018},
and essentially all of it lies below ${\sim}30$\,GHz. Above that frequency
the record is nearly empty. \citet{Steffes1994} and \citet{Mauersberger1996}
observed \SbPriorNStarA{} and \SbPriorNStarB{} nearby targets at
\SbPriorFreqGHz\,GHz, each in no more than \SbPriorBandMHz\,MHz of a single
line. The union of frequencies searched here is \SbPriorRatio{} times that
bandwidth, and the coverage of the median system alone, \SysUnionAMed\,GHz,
is still \SbPriorRatioSysMed{} times larger. Neither search allowed for drift or tested whether a signal
persisted. Of ALMA's receiver bands, only Band~3 has been searched, and that search
looked at distant stars in calibrator fields \citep{Mason2025}.

There are good reasons to look at these frequencies. Interstellar dispersion
and scintillation are negligible, and an aperture delivers far more gain per
watt. The costs are real as well, however. System temperatures are higher,
the atmosphere is opaque in places at both ends of the link, and the narrow
beam (\BmArcsec{} arcsec for the \BenchDiam\,m benchmark of
\S\ref{sec:mmdifferent}) would have to be aimed at the Earth deliberately.
Since nothing is known about such a transmitter, limits are stated here as
an equivalent isotropic radiated power (EIRP), which needs no model of it.
The targets were also chosen by other programmes for their own purposes, so
the sample cannot bound the fraction of stars that host a transmitter
\citep{Wright2018,Margot2023,EarthDetectingEarth2025}, as \S\ref{sec:bias}
shows.

The signal sought is an \emph{unresolved spectral carrier}: emission confined
to a few channels at the stellar position, present at more than one epoch,
and allowed to drift in frequency as it must for a transmitter on a rotating,
orbiting body (\S\ref{sec:statistic}). Broadband leakage on a planetary scale
\citep{BRaTs2026} lies outside the search by construction. An event is worth
pursuing only if it is found again \citep{Sheikh2021}, and the survey is
built around that test. Recurrence is not enough on its own, however, since
a molecular line or a source of interference can also return. \citet{Mason2025} identified the two
obstacles particular to this band: drift rates scale with frequency, and the
spectrum is crowded with molecular lines. The method therefore searches over
drift, compares every event with a line catalogue, and looks for the event
again in a separated epoch. As ALMA is an interferometer, the star can
also be compared with control positions inside its own primary beam, which a
single dish cannot resolve, and that comparison takes the place of the usual
interference rejection (\S\ref{sec:method}).

Experiments with different channel widths are compared here on the minimum
detectable received flux, $F_{\rm min}=S_{\rm min}\,\delta\nu$, each at its
own channel width. The comparison has to be made in this way because
$F_{\rm min}$ rises as $\delta\nu^{1/2}$, and a limit holds only for a signal
narrower than the channel that set it. The channels of the searches compared
span \LtChanSpanDex{} orders of magnitude, from \LtChanNarrow{} to
\LtChanWide, while this survey's own run from \SzChanFinestKHz\,kHz to
\SzChanWidestMHz\,MHz.

Figure~\ref{fig:context} draws each published threshold beside its value
rescaled to a common $\LtCommonSnr\sigma$. The rescaling is only an
arithmetic convention, since it does not match the false-alarm
probabilities of the different searches, and of the limits plotted only this
survey's $\mathrm{EIRP}_{90}$ is calibrated with injected carriers. The
plotted values are in \texttt{figures/figure\_data.csv} of the data archive
\citep{SIdeposit}.

The closest comparison is with the Band~3 search of \citet{Mason2025}, and
channel for channel this survey is the shallower of the two except in its
deepest windows. At a bare $\LtCommonSnr\sigma$ threshold that search reaches
a median \MasonFluxMed\,W\,m$^{-2}$ in a \MasonChanKHz\,kHz channel. The
deepest window here outside the line mask reaches \FluxTrigBest\,W\,m$^{-2}$
in \FluxChanBestKHz\,kHz, a factor \FluxNatBestRatio{} deeper, but the median
window reaches only \FluxTrigMed\,W\,m$^{-2}$ in \FluxChanMedKHz\,kHz, a
factor \FluxNatMedRatio{} shallower. What distinguishes this survey is
instead its frequency range, and its \NSysClassA{} systems within 40\,pc,
where the Band~3 search observed distant stars in \MasonNCal{} calibrator
fields over \MasonBandMHz\,MHz.
